% ============================================================
\chapter{Tests d'Hypoth\`eses Bay\'esiens}
\label{ch:tests}
% ============================================================

Le test d'hypoth\`ese fr\'equentiste repose sur la $p$-valeur~: la probabilit\'e d'observer des donn\'ees au moins aussi extr\^emes que celles obtenues, sous l'hypoth\`ese nulle. Mais cette quantit\'e r\'epond \`a une question que personne ne pose vraiment~: ce qu'on veut savoir, c'est la probabilit\'e que l'hypoth\`ese soit vraie \emph{\'etant donn\'e les donn\'ees}, pas l'inverse. Le test bay\'esien r\'epond directement \`a cette question gr\^ace au \emph{facteur de Bayes}~: le rapport des vraisemblances marginales sous chaque hypoth\`ese. Harold Jeffreys, dans son trait\'e de 1961, a syst\'ematis\'e cette approche et propos\'e l'\'echelle qui porte son nom pour interpr\'eter la force de l'\'evidence.

\begin{intuition}[title=\textbf{Id\'ee centrale}]
Le test bay\'esien compare les hypoth\`eses via leurs probabilit\'es a
posteriori ou le \textbf{facteur de Bayes}, qui mesure la force relative
des \'evidences en faveur de chaque hypoth\`ese, sans d\'ependre de la
r\`egle d'arr\^et ou de l'espace \'echantillonnal.
\end{intuition}

% ------------------------------------------------------------
\section{Formulation bay\'esienne du test}
% ------------------------------------------------------------

\begin{definition}[Test bay\'esien]
Soient deux hypoth\`eses $H_0 : \theta \in \Theta_0$ et
$H_1 : \theta \in \Theta_1$ avec $\Theta_0 \cup \Theta_1 = \Theta$
et $\Theta_0 \cap \Theta_1 = \emptyset$. On attribue des probabilit\'es
a priori $\PP(H_0) = \pi_0$ et $\PP(H_1) = \pi_1 = 1 - \pi_0$.
Les probabilit\'es a posteriori sont~:
\[
  \PP(H_k \mid x) = \frac{m_k(x)\,\pi_k}{m(x)},
  \quad k = 0, 1,
\]
o\`u $m_k(x) = \int_{\Theta_k} f(x \mid \theta)\,\pi(\theta \mid H_k)\,d\theta$.
\end{definition}

\begin{definition}[Facteur de Bayes]
Le \textbf{facteur de Bayes} en faveur de $H_0$ contre $H_1$ est~:
\[
  B_{01} = \frac{m_0(x)}{m_1(x)}
  = \frac{\PP(H_0 \mid x) / \PP(H_1 \mid x)}
         {\PP(H_0) / \PP(H_1)}.
\]
\end{definition}

\begin{formules}[title=\textbf{Rapport des cotes a posteriori}]
\[
  \underbrace{\frac{\PP(H_0 \mid x)}{\PP(H_1 \mid x)}}_{\text{cote post.}}
  = \underbrace{B_{01}}_{\text{facteur de Bayes}}
  \times
  \underbrace{\frac{\pi_0}{\pi_1}}_{\text{cote a priori}}.
\]
\end{formules}

% ------------------------------------------------------------
\section{Interpr\'etation du facteur de Bayes}
% ------------------------------------------------------------

\begin{formules}[title=\textbf{\'Echelle de Jeffreys pour $B_{01}$}]
\begin{center}
\renewcommand{\arraystretch}{1.3}
\begin{tabular}{cl}
\toprule
$\log_{10} B_{01}$ & \textbf{\'Evidence en faveur de $H_0$} \\
\midrule
$0$ \`a $1/2$ & Faible \\
$1/2$ \`a $1$ & Substantielle \\
$1$ \`a $2$ & Forte \\
$> 2$ & D\'ecisive \\
\bottomrule
\end{tabular}
\end{center}
\end{formules}

\begin{remarque}
Le facteur de Bayes est sym\'etrique~: $B_{10} = 1/B_{01}$.
Si $B_{01} > 1$, les donn\'ees soutiennent $H_0$~;
si $B_{01} < 1$, elles soutiennent $H_1$.
\end{remarque}

% ------------------------------------------------------------
\section{Test d'une hypoth\`ese ponctuelle}
% ------------------------------------------------------------

\begin{definition}[Hypoth\`ese ponctuelle]
On teste $H_0 : \theta = \theta_0$ contre $H_1 : \theta \neq \theta_0$.
On attribue une masse $\pi_0 = \PP(\theta = \theta_0) > 0$ et une
densit\'e continue $\pi_1(\theta)$ sur $\Theta_1$.
\end{definition}

\begin{theoreme}[Facteur de Bayes pour $H_0$ ponctuelle]
\[
  B_{01} = \frac{f(x \mid \theta_0)}
  {\int_{\Theta_1} f(x \mid \theta)\,\pi_1(\theta)\,d\theta}
  = \frac{L(\theta_0 \mid x)}{m_1(x)}.
\]
\end{theoreme}

\begin{attention}[title=\textbf{Paradoxe de Jeffreys--Lindley}]
Avec un prior vague ($\tau_0^2 \to \infty$) sous $H_1$, le facteur de
Bayes $B_{01} \to \infty$, favorisant toujours $H_0$, m\^eme quand le
$p$-value fr\'equentiste est tr\`es petit. Ce paradoxe illustre
l'importance du choix du prior sous $H_1$.
\end{attention}

\begin{exemple}
\textbf{Test $\mu = 0$ dans le mod\`ele gaussien.}
Soit $X_1, \ldots, X_n \overset{\text{iid}}{\sim} \mathcal{N}(\mu, \sigma^2)$
avec $\sigma^2$ connu. Sous $H_0 : \mu = 0$, sous $H_1 : \mu \sim \mathcal{N}(0, \tau^2)$.

Le facteur de Bayes est~:
\[
  B_{01} = \sqrt{\frac{\tau^2 + \sigma^2/n}{\sigma^2/n}}
  \exp\!\left(-\frac{n\bar{x}^2}{2\sigma^2}
  \cdot \frac{\tau^2}{\tau^2 + \sigma^2/n}\right).
\]
\end{exemple}

% ------------------------------------------------------------
\section{R\`egle de d\'ecision bay\'esienne}
% ------------------------------------------------------------

\begin{definition}[R\`egle de d\'ecision optimale]
Sous la perte 0--1~:
\[
  L(H_k, d) = \begin{cases}
  0 & \text{si } d = k, \\
  1 & \text{si } d \neq k,
  \end{cases}
\]
la d\'ecision optimale est~: accepter $H_0$ si
$\PP(H_0 \mid x) > \PP(H_1 \mid x)$, soit $B_{01} > \pi_1/\pi_0$.

Avec $\pi_0 = \pi_1 = 1/2$, cela revient \`a accepter $H_0$ si $B_{01} > 1$.
\end{definition}

\begin{theoreme}[R\`egle de d\'ecision sous perte asym\'etrique]
Sous la perte
$L(H_0, d=1) = c_0$ (erreur de type I) et $L(H_1, d=0) = c_1$ (erreur de type II),
la r\`egle optimale est~: accepter $H_0$ si
\[
  B_{01} > \frac{c_0}{c_1} \cdot \frac{\pi_1}{\pi_0}.
\]
\end{theoreme}

% ------------------------------------------------------------
\section{Calcul du facteur de Bayes}
% ------------------------------------------------------------

\subsection{Cas conjugu\'e}

\begin{proposition}[Facteur de Bayes en conjugaison]
Pour un mod\`ele conjugu\'e, la vraisemblance marginale est souvent
calculable analytiquement. Par exemple, pour le mod\`ele
Bernoulli--B\^eta avec $\theta \sim \text{Be}(a, b)$~:
\[
  m(x) = \frac{B(a + s, b + n - s)}{B(a, b)},
\]
o\`u $B(\cdot, \cdot)$ est la fonction b\^eta.
\end{proposition}

\subsection{M\'ethode de Savage--Dickey}

\begin{theoreme}[Rapport de densit\'e de Savage--Dickey]
Pour tester $H_0 : \theta = \theta_0$ contre
$H_1 : \theta \sim \pi_1(\theta)$, si $\pi_1$ est le prior sous $H_1$
et $\pi_1(\theta \mid x)$ le posterior correspondant~:
\[
  B_{01} = \frac{\pi_1(\theta_0 \mid x)}{\pi_1(\theta_0)}.
\]
\end{theoreme}

\begin{proof}
Par d\'efinition~:
\[
  B_{01} = \frac{f(x \mid \theta_0)}{m_1(x)}.
\]
Or $\pi_1(\theta_0 \mid x) = f(x \mid \theta_0)\pi_1(\theta_0)/m_1(x)$,
d'o\`u $f(x \mid \theta_0)/m_1(x) = \pi_1(\theta_0 \mid x)/\pi_1(\theta_0)$.
\end{proof}

\subsection{Approximations num\'eriques}

\begin{definition}[Approximation harmonique]
La vraisemblance marginale peut \^etre estim\'ee par la moyenne harmonique
des vraisemblances \'evalu\'ees sur les \'echantillons MCMC~:
\[
  \hat{m}(x) = \left(\frac{1}{T}\sum_{t=1}^T \frac{1}{L(\theta^{(t)} \mid x)}\right)^{-1}.
\]
\end{definition}

\begin{attention}[title=\textbf{Instabilit\'e de la moyenne harmonique}]
Cette estimation a une variance potentiellement infinie et ne doit
\^etre utilis\'ee qu'avec prudence. On pr\'ef\`ere des m\'ethodes plus
stables comme le \textit{bridge sampling} ou WAIC.
\end{attention}

% ------------------------------------------------------------
\section{Crit\`eres d'information bay\'esiens}
% ------------------------------------------------------------

\begin{definition}[BIC --- Bayesian Information Criterion]
\[
  \text{BIC} = -2\log L(\hat{\theta}_{\text{MLE}} \mid x) + k\log n,
\]
o\`u $k$ est le nombre de param\`etres. Le BIC est une approximation du
log-facteur de Bayes~: $-2\log B_{01} \approx \text{BIC}_0 - \text{BIC}_1$.
\end{definition}

\begin{definition}[WAIC --- Widely Applicable IC]
\[
  \text{WAIC} = -2\sum_{i=1}^n \log \E_{\theta \mid x}[f(x_i \mid \theta)]
  + 2\,p_{\text{WAIC}},
\]
o\`u $p_{\text{WAIC}} = \sum_{i=1}^n \Var_{\theta \mid x}[\log f(x_i \mid \theta)]$
est le nombre effectif de param\`etres.
\end{definition}

\begin{definition}[LOO-CV --- Leave-One-Out Cross-Validation]
\[
  \text{LOO} = -2\sum_{i=1}^n \log f(x_i \mid x_{-i}),
\]
o\`u $f(x_i \mid x_{-i}) = \int f(x_i \mid \theta)\,\pi(\theta \mid x_{-i})\,d\theta$
est la densit\'e pr\'edictive \textit{leave-one-out}.
\end{definition}

% ------------------------------------------------------------
\section{ROPE et \'equivalence pratique}
% ------------------------------------------------------------

\begin{definition}[ROPE --- Region of Practical Equivalence]
La \textbf{ROPE} est un intervalle $[\theta_0 - \varepsilon, \theta_0 + \varepsilon]$
d\'efinissant l'\'equivalence pratique \`a $\theta_0$. La d\'ecision
est~:
\begin{itemize}
  \item Accepter $H_0$ si $\text{HDI} \subset \text{ROPE}$.
  \item Rejeter $H_0$ si $\text{HDI} \cap \text{ROPE} = \emptyset$.
  \item Ind\'ecis sinon.
\end{itemize}
\end{definition}

% ------------------------------------------------------------
\section{Comparaison avec les $p$-valeurs}
% ------------------------------------------------------------

\begin{proposition}[Calibration Bayes--$p$-value]
Sellke, Bayarri et Berger (2001) montrent que~:
\[
  B_{01} \geq -\frac{1}{e \cdot p \cdot \log p}
  \quad \text{pour } p < 1/e,
\]
o\`u $p$ est le $p$-value. Ainsi, un $p$-value de $0.05$ correspond \`a
$B_{01} \geq 2.5$ au plus, soit une \'evidence faible \`a mod\'er\'ee
pour $H_0$.
\end{proposition}

% ------------------------------------------------------------
\section{Impl\'ementation Python}
% ------------------------------------------------------------

\begin{codeblock}[title=\textbf{Facteur de Bayes pour un test gaussien}]
\begin{minted}{python}
import numpy as np
from scipy import stats

# Données
np.random.seed(42)
n = 25
sigma = 1.0
mu_true = 0.3
data = np.random.normal(mu_true, sigma, n)
x_bar = data.mean()

# H0: mu = 0, H1: mu ~ N(0, tau^2)
tau = 1.0

# Facteur de Bayes analytique
var_ratio = tau**2 / (tau**2 + sigma**2/n)
log_BF01 = 0.5 * np.log(1 + n*tau**2/sigma**2) \
           - 0.5 * n * x_bar**2 / sigma**2 * var_ratio
BF01 = np.exp(log_BF01)

print(f"x_bar = {x_bar:.4f}")
print(f"BF01 = {BF01:.4f}")
print(f"log10(BF01) = {np.log10(BF01):.4f}")

if BF01 > 1:
    print("Données en faveur de H0")
else:
    print(f"Données en faveur de H1 (BF10 = {1/BF01:.4f})")

# Comparaison p-value
t_stat = x_bar / (sigma / np.sqrt(n))
p_value = 2 * (1 - stats.norm.cdf(abs(t_stat)))
print(f"p-value (bilatéral): {p_value:.4f}")
\end{minted}
\end{codeblock}

\begin{codeblock}[title=\textbf{ROPE et HDI avec PyMC}]
\begin{minted}{python}
import pymc as pm
import arviz as az
import matplotlib.pyplot as plt

# Modèle bayésien
with pm.Model() as model_test:
    mu = pm.Normal("mu", mu=0, sigma=10)
    y = pm.Normal("y", mu=mu, sigma=sigma,
                  observed=data)
    trace = pm.sample(3000, random_seed=42)

# HDI et ROPE
hdi = az.hdi(trace, var_names=["mu"], hdi_prob=0.95)
print("95% HDI for mu:", hdi)

# ROPE test
rope = [-0.1, 0.1]
az.plot_posterior(trace, var_names=["mu"],
                 rope=rope, ref_val=0,
                 figsize=(8, 4))
plt.savefig("rope_test.pdf")

# Proportion du posterior dans la ROPE
mu_samples = trace.posterior["mu"].values.flatten()
in_rope = np.mean((mu_samples >= rope[0]) &
                   (mu_samples <= rope[1]))
print(f"P(mu in ROPE | data) = {in_rope:.4f}")
\end{minted}
\end{codeblock}

\begin{codeblock}[title=\textbf{Comparaison de mod\`eles avec WAIC}]
\begin{minted}{python}
# Deux modèles en compétition
with pm.Model() as model_linear:
    a = pm.Normal("a", 0, 10)
    b = pm.Normal("b", 0, 10)
    sigma_m = pm.HalfNormal("sigma", 5)
    x_pred = np.linspace(0, 10, n)
    mu_pred = a + b * x_pred
    y_obs = pm.Normal("y", mu=mu_pred, sigma=sigma_m,
                      observed=data)
    trace_lin = pm.sample(2000, random_seed=42)

with pm.Model() as model_null:
    a = pm.Normal("a", 0, 10)
    sigma_m = pm.HalfNormal("sigma", 5)
    y_obs = pm.Normal("y", mu=a, sigma=sigma_m,
                      observed=data)
    trace_null = pm.sample(2000, random_seed=42)

# Comparaison WAIC
waic_lin = az.waic(trace_lin, model_linear)
waic_null = az.waic(trace_null, model_null)
comp = az.compare({"linear": trace_lin,
                    "null": trace_null})
print(comp)
\end{minted}
\end{codeblock}

% ------------------------------------------------------------
\section{Exercices}
% ------------------------------------------------------------

\begin{exercice}
Calculer le facteur de Bayes pour le test $H_0 : \theta = 1/2$
contre $H_1 : \theta \sim \text{Be}(1,1)$ dans le mod\`ele
Bernoulli avec $n = 20$ et $s = 14$.
\end{exercice}

\begin{exercice}
D\'emontrer le rapport de Savage--Dickey. V\'erifier num\'eriquement
sur le mod\`ele Normal--Normal.
\end{exercice}

\begin{exercice}
Illustrer le paradoxe de Jeffreys--Lindley~: pour $n = 100$,
$\bar{x} = 0.2$, $\sigma = 1$, calculer le $p$-value et le facteur de
Bayes $B_{01}$ pour diff\'erentes valeurs de $\tau \in \{1, 10, 100\}$.
Commenter.
\end{exercice}

\begin{exercice}
\textbf{(Num\'erique)} Comparer WAIC et LOO-CV pour la s\'election
entre un mod\`ele lin\'eaire et un mod\`ele quadratique sur des donn\'ees
simul\'ees.
\end{exercice}

\begin{exercice}
D\'eriver l'approximation BIC du facteur de Bayes \`a partir de
l'approximation de Laplace de la vraisemblance marginale.
\end{exercice}
